\section{离散动力系统、分岔与混沌}
\label{sec:15-Differential-Equations-and-Dynamical-Systems/02-Stability-and-Dynamical-Systems/03}

\subsection{把时间从连续拧成离散，世界变大了}

连续时间的变化规则是微分方程 \(\dot u = F(u)\)。上一单元用 Logistic 方程描述种群增长：
\[
\dot x = rx\left(1-\frac{x}{K}\right),
\]
任何正的初值都单调收敛到承载力 \(K\)。从任何起点出发，轨道都平静地滑向同一个归宿——无聊，但可预测。

现在做一个很自然的改变。许多物种的繁殖是季节性的、不重叠代际的：今年的种群数量决定明年的。把连续方程中的瞬时变化率换成逐年更新：
\[
x_{n+1} = f(x_n).
\]
取和 Logistic 方程同样的建模思想——增长受资源限制，写成离散版就是 Logistic 映射：
\[
x_{n+1} = r\,x_n(1-x_n), \qquad x_n\in[0,1],\quad 0\le r\le 4,
\]
其中 \(x_n\) 是第 \(n\) 代种群占环境上限的比例，参数 \(r\) 综合了繁殖率（\(r>4\) 时轨道会跑出 \([0,1]\)，模型失效）。

数学结构上看，这不过是最简单的非线性迭代——单个变量、单个参数、一个二次函数。连续系统中单个变量的一维自治流，解曲线互不相交，直线上无处可绕，轨道只能单调走向平衡点。离散迭代没有这把锁。下面你会看到，同一个建模思想，只需要拧一个参数 \(r\)，就能上演连续一维系统想都不敢想的全部剧情。

\subsection{不动点和乘子：离散版的稳定判据}

不动点满足 \(x^* = f(x^*)\)——迭代到这里就停了，它是连续系统平衡点 \(F(u^*)=0\) 的离散对应物。

判断一个不动点是吸引还是排斥邻近轨道，在线性化。令 \(e_n = x_n - x^*\)，由 \(f(x^*) = x^*\) 和 Taylor 展开：
\[
e_{n+1} = f(x^*+e_n) - f(x^*) = f'(x^*)\,e_n + o(|e_n|).
\]
小扰动每一步大约被乘以 \(f'(x^*)\)——这个数叫不动点的乘子（multiplier）。

判据是干净的：
\begin{itemize}
\tightlist
\item \(|\lambda| < 1\)：扰动每步几何衰减，不动点吸引邻近轨道；
\item \(|\lambda| > 1\)：扰动每步几何放大，不动点排斥邻近轨道；
\item \(|\lambda| = 1\)：临界情形，线性化不提供结论——与连续情形``特征值实部为零''完全平行。
\end{itemize}

比较一下连续和离散的口诀：连续看特征值实部，实部为负则稳定；离散看乘子的模，模小于 1 则稳定。离散多出一种连续一维不可能有的现象：乘子为负时，\(e_{n+1} \approx -|\lambda| e_n\)，扰动每步翻转符号——轨道以振荡的方式靠近或远离不动点。迭代允许``走过头然后被拉回来''，而一维连续流不允许掉头（你没法在直线上 U 型转弯而不撞上自己）。

\subsection{Logistic 映射：一个参数拧出的全部剧情}

对 Logistic 映射 \(f(x)=rx(1-x)\)，先找不动点：
\[
x^* = r x^*(1-x^*) \implies x^*=0,\quad\text{以及 } r>1 \text{ 时的 } x^*=1-\frac{1}{r}.
\]
乘子 \(f'(x) = r(1-2x)\)，逐段查看：

\textbf{\(0\le r < 1\)：灭绝。} \(x^*=0\) 是 \([0,1]\) 内唯一不动点，乘子绝对值为 \(r<1\)。任何初值逐代衰减到零——种群无法维持。

\textbf{\(1 < r < 3\)：收敛到正定态。} \(x^*=1-1/r\) 成为主角。乘子 \(f'(x^*) = r(1-2(1-1/r)) = 2-r\)，其绝对值在 \(r\in(1,3)\) 内小于 1。有趣的是 \(r>2\) 时乘子已经为负——种群以振荡方式逼近定态，每一代在定态两侧来回，幅度逐步收窄。连续 Logistic 方程从来不会振荡，它是单调饱和的；离散版从 \(r=2\) 起就出现了连续版没有的符号交替。

\textbf{\(r=3\)：第一次失稳。} 乘子恰好为 \(-1\)，定态失去吸引力。之后的剧情连续一维系统永远演不出。

\textbf{\(3 < r < 1+\sqrt{6}\approx 3.449\)：2-周期。} 轨迹在两个值之间稳定地来回跳动。单次迭代的不动点失稳了，但复合两次的映射 \(f\circ f\) 生出了一对新的稳定不动点。这对点的稳定判据仍看乘子——链式法则给出 \((f\circ f)'(x_1) = f'(x_1)f'(x_2)\)，计算表明其模在此区间内小于 1。

\textbf{\(r\) 继续增大：倍周期级联。} 2-周期失稳，4-周期接班；4-周期失稳，8-周期接班；8 变 16，16 变 32……每一次接棒，都是复合映射的乘子再次穿过 \(-1\)。这串倍周期分岔的参数值 \(r_1, r_2, r_3, \ldots\) 越来越密集。

密集的节奏不是随意的。相邻分岔间距之比趋向一个常数：
\[
\delta = \lim_{k\to\infty} \frac{r_k - r_{k-1}}{r_{k+1} - r_k} \approx 4.669.
\]
它就是 Feigenbaum 常数。它不专属于 Logistic 映射——任何带一个光滑单峰的单参数映射族，只要峰是二次的（即 \(f''\) 在峰处不为零），都给出一模一样的 \(\delta\)。这不是巧合，是深刻的普适性（universality）：倍周期分岔的``时间表''由映射在峰附近的函数形状所决定，而很多映射的峰的形状在重参数化意义下是同一类。

\textbf{\(r > r_\infty \approx 3.570\)：混沌区域。} 累积点之后，多数的参数值上轨道永不重复、永不安定。它在区间中看似随机地游荡，但每一个下一步都被 \(f\) 完全确定。这里没有噪声源，没有外界扰动——确定的规则、确定的初值，行为却像随机过程。其中还嵌着狭窄的周期窗口：比如 \(r\approx 3.83\) 附近突然出现一段稳定的 3-周期，但参数稍挪，3-周期又失稳并再次通向混沌。

\subsection{Lyapunov 指数：把``对初值敏感''变成数}

混沌的一个标志性特征是对初值极端敏感。两条初值相差 \(\varepsilon\) 的轨道，迭代 \(n\) 步后的分离程度由复合映射的导数决定：
\[
|f^n(x_0+\varepsilon) - f^n(x_0)|
\approx |\varepsilon| \cdot \bigl|(f^n)'(x_0)\bigr|.
\]
由链式法则：
\[
(f^n)'(x_0) = \prod_{k=0}^{n-1} f'(x_k),
\]
其中 \(x_k = f^k(x_0)\) 是轨道上的点。分离率是沿途各点导数的累积乘积。取对数化为时间平均：
\[
\lambda = \lim_{n\to\infty} \frac{1}{n}\sum_{k=0}^{n-1} \ln|f'(x_k)|,
\]
这就是沿该轨道的 Lyapunov 指数。

\(\lambda > 0\) 意味着典型的小扰动沿某些方向指数放大——混沌的重要诊断。稳定周期轨道对应的乘子模小于 1，Lyapunov 指数为负，邻近轨道被吸拢。对 Logistic 映射在 \(r=4\) 时，其自然不变测度下的典型轨道 Lyapunov 指数为 \(\lambda = \ln 2 \approx 0.693\)。

翻译成预测的时间尺度：假设初值只测准到 \(10^{-10}\) 的相对精度，而 \(\lambda = \ln 2\) 意味着每步分离因子约为 2。叠 30 余步后，由初值不确定性造成的误差就可能胀到 \(O(1)\)——轨道级别的预测有效期极为有限。

这里需要划清一条线。如果初值和规则都能以无限精度给出，确定性轨道当然仍然被唯一决定。问题不是计算能力，是信息——现实中的初值永远只有有限精度，而在正 Lyapunov 指数下，这个精度很快就被耗尽。能算出准确的下一步（如果给你精确的当前值）和能预言 30 步后的状态，是两件完全不同的事。

因此研究混沌系统时常转向统计量：不变测度、访问频率、分形维数。这些量在轨道细节不可预测时，仍然可能是稳定的。但统计描述需要遍历性等额外条件——不是所有混沌系统都方便地遍历。

\subsection{Poincaré 截面：把连续流压扁成离散映射}

离散工具的价值不止于分析本来就在离散时间中的系统。对连续流（三维以上自主系统），在相空间中搁一张截面，每当轨道沿同一方向穿过它就记录一个交点：连续时间流被压缩成截面上的离散映射——Poincaré 首次回归映射。

极限环在截面上的交点是该映射的不动点，稳定性仍由乘子模判定。本节关于不动点、分岔与混沌的全部词汇由此可以搬回连续系统。三维及以上连续系统中的混沌——Lorenz 吸引子是最著名的例子——正是通过 Poincaré 截面被解开的：看截面上的点如何散布、拉伸、折叠，用离散工具理解连续混沌。

\subsection{落点：可计算、可预测和可控制是三件不同的事}

本节串联起来的逻辑链值得单独记住：一个规则完全确定的系统，下一步总是可计算的；但 Lyapunov 指数为正时，长期轨道不可预测；而不可预测的系统，仍然有可能通过反馈控制在统计意义上被稳定——控制和预测不是一套问题。

回过头看上一单元末尾的伏笔：梯度下降 \(u_{k+1} = u_k - \alpha \nabla f(u_k)\) 正是离散动力系统，极小点是它的不动点。一维情况下，乘子为 \(1 - \alpha f''(u^*)\)，稳定条件 \(|\lambda|<1\) 即 \(0 < \alpha < 2/f''(u^*)\)。当步长越过这个上限时乘子穿过 \(-1\)，不动点经倍周期分岔失稳——梯度下降在极小点附近来回蹦跳而不是收敛。这是\hyperref[sec:09-Convex-Optimization/06-Optimization-Algorithms/01]{``梯度下降真正难在步长''}的最深层原因：步长不当不只是慢，而是会触发离散化制造的动力学分岔。同一个判据也解释过数值分析中的现象——显式 Euler 用于 \(y'=-y\) 时，步长 \(h>2\) 时数值解反而增长：乘子 \(1-h\) 越过了 \(-1\)（见\hyperref[ch:066]{``常微分方程数值方法：沿局部斜率走完整条轨迹''}）。

Logistic 映射教给我们最大的教训在这里：连续和离散的门槛不是计算细节，是动力学丰富度的分界线。描述同一个种群增长的想法，在连续时间中只有单调收敛一种结局；在离散时间中却能上演定态、周期、倍周期分岔和混沌的完整剧目。一切只因为你把``每时每刻都在变化''换成了``一年变一次''。
